\documentclass[11pt, a4paper]{article}

\input{bfw-style.tex}

\title{\vspace*{-1.5cm}\Huge\bfseries\color{bfwPrimaryDark}Aufgabenheft Session~3\\[0.4em]
  \large\color{bfwInkSoft}\textnormal{Betrieb, Unternehmen \& Unternehmensziele --- Übungen im IHK-Klausurformat}}
\author{\textsf{IT Lernfeld 1 --- Das Unternehmen und die eigene Rolle im Betrieb beschreiben}}
\date{Selbstlernmaterial}

\begin{document}
\maketitle
\thispagestyle{empty}

\begin{merksatz}[title={\bfseries So nutzen Sie dieses Aufgabenheft}]
Bearbeiten Sie alle Aufgaben zunächst \emph{ohne} in die Lösungen zu schauen. Schreiben
Sie Ihre Antworten in vollständigen Sätzen --- die IHK-Prüfung verlangt formulierte
Begründungen, nicht bloße Stichworte. Beachten Sie die \emph{Operatoren} (nennen,
erläutern, beschreiben, begründen, beurteilen) --- sie geben den geforderten
Antwort-Tiefegrad vor. Erst danach öffnen Sie den Lösungsteil ab Seite~\pageref{loesungen}
und vergleichen.
\end{merksatz}

\bigskip

\textbf{Kontext:} Sie sind seit einigen Wochen als neue Auszubildende bei der \textbf{JIKU IT-Solutions GmbH},
einem IT-Systemhaus mit 52 Mitarbeitenden am Hamburger Standort, davon 8 Auszubildende, tätig.
Ihre Ausbilderin, Frau Meyerhoff, gibt Ihnen heute einen Überblick über den Betrieb und seine Ziele.

\bigskip

\noindent
\begin{tabular}{|l|l|l|l|l|l|l|l|l|}
\hline
\textbf{Aufgabe} & 1 & 2 & 3 & 4 & 5 & 6 & 7 & \textbf{Gesamt}\\
\hline
\textbf{Punkte} & /10 & /8 & /8 & /10 & /6 & /8 & /6 & /56\\
\hline
\end{tabular}

\tableofcontents
\vspace{1em}
\hrule

\section{Aufgaben}

\subsection*{Aufgabe~1 --- JIKU IT-Solutions kennenlernen (10 Punkte)}

\begin{enumerate}[label=(\alph*)]
  \item \textbf{Nennen} Sie vier Eckdaten des JIKU-Systemhausverbundes: Anzahl der Systemhäuser,
  Standorte, Gesamtmitarbeiterzahl und Gesamtzahl der Auszubildenden. \hfill (4~P.)

  \vspace{3.5cm}

  \item \textbf{Erläutern} Sie, was ein IT-Systemhaus von einem reinen IT-Hersteller unterscheidet,
  und nennen Sie zwei Leistungsbereiche aus dem Portfolio von JIKU IT-Solutions. \hfill (4~P.)

  \vspace{4cm}

  \item \textbf{Beschreiben} Sie, was bei JIKU unter Corporate Social Responsibility (CSR) zu
  verstehen ist. Nennen Sie zwei konkrete Maßnahmen. \hfill (2~P.)

  \vspace{3cm}
\end{enumerate}

\subsection*{Aufgabe~2 --- Begriffsabgrenzung (8 Punkte)}

Kreuzen Sie an, ob folgende Aussagen richtig ({[~]} R) oder falsch ({[~]} F) sind.
Korrigieren Sie falsche Aussagen in einem Satz.

\bigskip

\noindent
\begin{tabular}{|p{8.5cm}|c|c|p{4cm}|}
\hline
\textbf{Aussage} & \textbf{R} & \textbf{F} & \textbf{Korrektur (wenn falsch)}\\
\hline
a) Ein Betrieb ist die rechtlich-wirtschaftliche Einheit, die Steuern zahlt. & {[~]} & {[~]} & \\[0.8cm]
\hline
b) Ein Unternehmen kann mehrere Betriebe haben. & {[~]} & {[~]} & \\[0.8cm]
\hline
c) Die Firma ist nach § 17 HGB der Name, unter dem ein Kaufmann seine Geschäfte betreibt. & {[~]} & {[~]} & \\[0.8cm]
\hline
d) JIKU IT-Solutions GmbH ist ein gemeinwirtschaftliches Unternehmen. & {[~]} & {[~]} & \\[0.8cm]
\hline
e) IT-Systemhäuser gehören dem tertiären Wirtschaftssektor an. & {[~]} & {[~]} & \\[0.8cm]
\hline
f) Erwerbswirtschaftliche Unternehmen dürfen keinen Gewinn anstreben. & {[~]} & {[~]} & \\[0.8cm]
\hline
g) Ein Konzern besteht aus mehreren rechtlich selbstständigen Unternehmen unter einer gemeinsamen Leitung. & {[~]} & {[~]} & \\[0.8cm]
\hline
h) Kartelle sind generell erlaubt, wenn alle beteiligten Unternehmen zustimmen. & {[~]} & {[~]} & \\[0.8cm]
\hline
\end{tabular}

\bigskip

\subsection*{Aufgabe~3 --- Wirtschaftliche Verflechtungen (8 Punkte)}

Frau Meyerhoff erklärt, dass JIKU IT-Solutions überlegt, einen kleineren Mitbewerber zu übernehmen
und mit einem Rechenzentrum-Anbieter ein befristetes Gemeinschaftsprojekt zur Cloud-Standardisierung
durchzuführen.

\begin{enumerate}[label=(\alph*)]
  \item \textbf{Ordnen} Sie beide Vorhaben den passenden Formen wirtschaftlicher Verflechtungen
  (Fusion, Konzern, Kartell, Konsortium, Franchise) zu und \textbf{begründen} Sie Ihre Zuordnung. \hfill (4~P.)

  \vspace{5cm}

  \item \textbf{Erläutern} Sie den Unterschied zwischen einer horizontalen und einer vertikalen
  Unternehmensverflechtung. Geben Sie je ein Beispiel aus der IT-Branche. \hfill (4~P.)

  \vspace{5cm}
\end{enumerate}

\subsection*{Aufgabe~4 --- Leitbild und Unternehmensziele (10 Punkte)}

\begin{enumerate}[label=(\alph*)]
  \item \textbf{Nennen} Sie die drei Hauptbestandteile eines Unternehmensleitbildes und
  \textbf{erläutern} Sie deren Unterschied an einem selbst gewählten Beispiel. \hfill (6~P.)

  \vspace{6cm}

  \item Was bedeutet das Akronym SMART in der Zielformulierung? \textbf{Formulieren} Sie
  anschließend ein SMART-Ziel für JIKU IT-Solutions. \hfill (4~P.)

  \vspace{5cm}
\end{enumerate}

\subsection*{Aufgabe~5 --- Sachziele und Formalziele (6 Punkte)}

Ordnen Sie die folgenden Unternehmensziele den Kategorien \textbf{Sachziel} (S) oder
\textbf{Formalziel} (F) zu und geben Sie an, ob es sich um ein \textbf{strategisches} (St)
oder \textbf{operatives} (O) Ziel handelt.

\bigskip

\noindent
\begin{tabular}{|p{9cm}|c|c|}
\hline
\textbf{Ziel} & \textbf{S oder F} & \textbf{St oder O}\\
\hline
a) JIKU möchte in diesem Quartal 50 neue Help-Desk-Verträge abschließen. & & \\[0.5cm]
\hline
b) JIKU strebt langfristig eine Eigenkapitalrentabilität von 15 \% an. & & \\[0.5cm]
\hline
c) Die Anzahl der Systemausfälle bei betreuten Kunden soll um 20 \% gesenkt werden. & & \\[0.5cm]
\hline
d) JIKU will in den nächsten 5 Jahren Marktführer im Hamburger IT-Systemhaus-Segment werden. & & \\[0.5cm]
\hline
e) Die Reaktionszeit des Help-Desks soll von 4 auf 2 Stunden halbiert werden. & & \\[0.5cm]
\hline
f) Der monatliche Umsatz soll im nächsten Quartal um 8 \% gesteigert werden. & & \\[0.5cm]
\hline
\end{tabular}

\bigskip

\subsection*{Aufgabe~6 --- Zielbeziehungen (8 Punkte)}

Bestimmen Sie für folgende Zielpaare die Zielbeziehung (\textbf{komplementär},
\textbf{konkurrierend} oder \textbf{indifferent}) und \textbf{begründen} Sie Ihre
Antwort in einem Satz. \hfill (je 2~P.)

\begin{enumerate}[label=(\alph*)]
  \item JIKU möchte die Anzahl der Auszubildenden erhöhen \emph{und} gleichzeitig die
  Ausbildungskosten senken.

  \vspace{2.5cm}

  \item JIKU plant höhere Investitionen in IT-Sicherheit \emph{und} möchte gleichzeitig
  die Kundenzufriedenheit steigern.

  \vspace{2.5cm}

  \item JIKU möchte mehr Innovationsprojekte starten \emph{und} gleichzeitig den
  Personalaufwand reduzieren.

  \vspace{2.5cm}

  \item JIKU möchte die Mitarbeiterzufriedenheit verbessern \emph{und} gleichzeitig
  die Produktivität steigern.

  \vspace{2.5cm}
\end{enumerate}

\subsection*{Aufgabe~7 --- Wirtschaftliche Kennzahlen (6 Punkte)}

Frau Meyerhoff teilt folgende Zahlen für den Hamburger JIKU-Standort mit:
\begin{itemize}
  \item Erträge im letzten Monat: 180.000~€
  \item Aufwendungen im letzten Monat: 150.000~€
  \item Jahresüberschuss: 36.000~€
  \item Eingesetztes Kapital: 240.000~€
  \item Support-Tickets: 600 in 8 Stunden (gestern); 540 in 6 Stunden (vorgestern)
\end{itemize}

\begin{enumerate}[label=(\alph*)]
  \item \textbf{Berechnen} Sie die Wirtschaftlichkeit des Hamburger Standorts für den
  letzten Monat. \hfill (2~P.)

  \vspace{3.5cm}

  \item \textbf{Berechnen} Sie die Rentabilität des eingesetzten Kapitals für das letzte
  Jahr. \hfill (2~P.)

  \vspace{3.5cm}

  \item \textbf{Berechnen} Sie die Produktivität des Support-Teams für beide Tage und
  \textbf{bestimmen} Sie, an welchem Tag das Team produktiver war. \hfill (2~P.)

  \vspace{4cm}
\end{enumerate}

% =====================================================================
\newpage
\section*{Lösungshinweise für Lehrende}
\label{loesungen}
\addcontentsline{toc}{section}{Lösungshinweise für Lehrende}
% =====================================================================

\begin{merksatz}[title={\bfseries Hinweis zur Nutzung}]
Dieser Abschnitt ist ausschließlich für Lehrende bestimmt. Die Lösungen geben
Musterlösungen und typische Fehlerquellen an. Teilnehmende erhalten diesen Teil
erst nach der eigenständigen Bearbeitung.
\end{merksatz}

\subsection*{Lösung Aufgabe~1}

\begin{enumerate}[label=(\alph*)]
  \item Vier Eckdaten des JIKU-Verbundes: \textbf{10 Systemhäuser}, \textbf{16 Standorte} in
  Deutschland, \textbf{480 Mitarbeiterinnen und Mitarbeiter} insgesamt,
  \textbf{96 Auszubildende} im Verbund.

  \textit{Häufiger Fehler: Mitarbeiterzahl am Hamburger Standort (52) wird mit Gesamtzahl (480) verwechselt.}

  \item Ein IT-Systemhaus agiert als \textbf{Mittler zwischen IT-Herstellern und Endkunden}: Es berät,
  beschafft Produkte und richtet diese ein. Ein IT-Hersteller \emph{produziert} dagegen
  Standardprodukte. Zwei Leistungsbereiche (Auswahl): IT-Service, Cloud-Hosting, Server \& Storage,
  IT-Infrastrukturen, IT-Beratung.

  \item CSR bezeichnet bei JIKU die \textbf{gelebte Wertschätzung für verantwortungsvolles und
  nachhaltiges Handeln}. Zwei Maßnahmen (Auswahl): flexible Arbeitszeiten / Home-Office,
  Betriebssport, Girls Day, Hackathon-Sponsoring, Stipendien, Fahrradleasing.
\end{enumerate}

\subsection*{Lösung Aufgabe~2}

\begin{itemize}
  \item a) \textbf{F} --- Der Betrieb ist die \emph{organisatorisch-technische} Einheit; das Unternehmen ist die rechtlich-wirtschaftliche.
  \item b) \textbf{R}
  \item c) \textbf{R}
  \item d) \textbf{F} --- JIKU ist ein \emph{erwerbswirtschaftliches} Unternehmen, das Gewinn anstrebt.
  \item e) \textbf{R}
  \item f) \textbf{F} --- Erwerbswirtschaftliche Unternehmen \emph{streben ausdrücklich nach Gewinn}.
  \item g) \textbf{R}
  \item h) \textbf{F} --- Preis- und Gebietskartelle sind nach dem Gesetz gegen Wettbewerbsbeschränkungen (GWB) \emph{verboten}, unabhängig von der Zustimmung der Beteiligten.
\end{itemize}

\subsection*{Lösung Aufgabe~3}

\begin{enumerate}[label=(\alph*)]
  \item Übernahme des Mitbewerbers: \textbf{Fusion} (Zusammenschluss zu einer rechtlichen Einheit)
  oder Begründung als Konzern möglich, wenn rechtliche Selbstständigkeit erhalten bleibt.
  Befristetes Cloud-Standardisierungsprojekt: \textbf{Konsortium} (vorübergehender Zusammenschluss
  für ein gemeinsames Vorhaben). \textit{Kartell oder Franchise wären hier falsch.}

  \item \textbf{Horizontal}: Unternehmen der gleichen Leistungsstufe (z.\,B.\ zwei Rechenzentren
  fusionieren). \textbf{Vertikal}: Unternehmen verschiedener Leistungsstufen (z.\,B.\ ein
  Chip-Hersteller übernimmt einen Server-Systemintegrator).
\end{enumerate}

\subsection*{Lösung Aufgabe~4}

\begin{enumerate}[label=(\alph*)]
  \item \textbf{Vision}: Wohin soll sich das Unternehmen entwickeln? (Zukunftsbild, z.\,B.\ „Marktführer in …")\\
  \textbf{Mission}: Was macht das Unternehmen und warum? (Zweck und Kernkompetenzen)\\
  \textbf{Werte}: Welche Grundsätze und Verhaltensweisen leiten das Handeln?\\
  Beispiel JIKU: Vision = „bestes IT-Systemhaus der Region", Mission = „wirtschaftlich, kompetent,
  sozial nachhaltig und umweltfreundlich handeln", Werte = Kundenorientierung, Exzellenz, Nachhaltigkeit.

  \item SMART: \textbf{S}pezifisch, \textbf{M}essbar, \textbf{A}usführbar/Akzeptiert, \textbf{R}ealistisch,
  \textbf{T}erminierbar. Beispiel-SMART-Ziel: „JIKU Hamburg erhöht die Anzahl betreuter
  IT-Systeme bis zum 31.~Dezember des laufenden Jahres von 400 auf 480 (20\,\% Steigerung),
  gemessen an abgeschlossenen Service-Level-Agreements."
\end{enumerate}

\subsection*{Lösung Aufgabe~5}

\begin{itemize}
  \item a) S / O (konkrete Leistung kurzfristig)
  \item b) F / St (wirtschaftlicher Erfolg langfristig)
  \item c) S / O (Qualitätsziel kurzfristig)
  \item d) F / St (Marktposition langfristig)
  \item e) S / O (Leistungsqualität kurzfristig)
  \item f) F / O (Umsatz kurzfristig)
\end{itemize}

\textit{Hinweis: Die Einordnung S/F und St/O ist manchmal kontextabhängig; gut begründete Abweichungen
können anerkannt werden.}

\subsection*{Lösung Aufgabe~6}

\begin{enumerate}[label=(\alph*)]
  \item \textbf{Konkurrierend} --- Mehr Auszubildende bedeutet höhere Ausbildungskosten; die Ziele
  behindern sich.
  \item \textbf{Komplementär} --- Bessere IT-Sicherheit reduziert Systemausfälle und Datenverluste,
  was die Kundenzufriedenheit steigert.
  \item \textbf{Konkurrierend} --- Innovationsprojekte erfordern Fachkräfte; Personalabbau
  entzieht diesen Projekten die nötige Kapazität.
  \item \textbf{Komplementär} --- Zufriedene Mitarbeiter sind motivierter und arbeiten produktiver;
  die Ziele verstärken sich.
\end{enumerate}

\subsection*{Lösung Aufgabe~7}

\begin{enumerate}[label=(\alph*)]
  \item Wirtschaftlichkeit $= \dfrac{180.000}{150.000} = \mathbf{1{,}2}$ (bzw.\ 120\,\%).
  Der Standort arbeitet wirtschaftlich (Wert $>1$).

  \item Rentabilität $= \dfrac{36.000 \times 100}{240.000} = \mathbf{15\,\%}$.

  \item Produktivität gestern $= \dfrac{600}{8} = \mathbf{75}$ Tickets/Stunde.\\
  Produktivität vorgestern $= \dfrac{540}{6} = \mathbf{90}$ Tickets/Stunde.\\
  Das Team war \textbf{vorgestern} produktiver.
\end{enumerate}

\end{document}
